% ============================================================
\section{Exceptions \& Smart Pointer (Abschnitt VII)}
% ============================================================
In C musste nach fast jedem Aufruf der R\"uckgabewert gepr\"uft werden -- Normal- und Fehlerfall vermischen sich, Kontext fehlt. C\texttt{++} trennt beides mit \textbf{Exceptions}.

\subsection{Exception Handling: \texttt{throw} / \texttt{try} / \texttt{catch}}
\begin{definition}[: Exception]
Eine Exception ist ein \textbf{Objekt}, das eine Fehlersituation repr\"asentiert und den normalen Kontrollfluss unterbricht, um an geeigneter Stelle behandelt zu werden.
\end{definition}
Ablauf in drei Schritten:
\begin{itemize}
  \item \code{try \{ ... \}} -- markiert einen Block, in dem ein Fehler auftreten kann.
  \item \code{throw obj;} -- wirft eine Exception.
  \item \code{catch (Typ e) \{ ... \}} -- f\"angt sie und behandelt sie.
\end{itemize}
\begin{lstlisting}
int divide(int a, int b) {
    if (b == 0) throw std::invalid_argument("Division durch 0");
    return a / b;
}
int main() {
    try {
        std::cout << divide(10, 0) << '\n';
    } catch (const std::exception& error) {   // const Referenz!
        std::cout << "Fehler: " << error.what();
    }
    return 0;   // Programm laeuft nach dem catch normal weiter
}
\end{lstlisting}
\begin{merke}
Wird eine Exception nicht lokal behandelt, wandert sie den Aufrufstack nach oben (\emph{Stack Unwinding}); dabei werden lokale Objekte korrekt \"uber ihre Destruktoren zerst\"ort. Wird sie \textbf{gar nicht} gefangen, beendet \code{std::terminate()} das Programm. Auf jeden \code{try}-Block muss direkt ein \code{catch}-Block folgen.
\end{merke}

\subsubsection*{Mehrere catch-Bl\"ocke}
Werden von oben nach unten gepr\"uft; der erste passende wird ausgef\"uhrt, der Rest \"ubersprungen. Ein allgemeiner Auffangblock steht \textbf{immer am Ende}.
\begin{lstlisting}
try { /* ... */ }
catch (const std::invalid_argument& e) { /* spezifisch */ }
catch (const std::exception& e)         { /* allgemeiner, bevorzugt */ }
catch (...)                             { /* letzte Sicherheitsstufe */ }
\end{lstlisting}

\subsection{Standard-Exceptions}
Einfache Typen (\code{throw 42;}) sind ungeeignet: keine Aussagekraft, keine Zusatzinfos, schlechte Unterscheidung. Besser: Klassen aus \code{<stdexcept>}, die von \code{std::exception} erben und die virtuelle Methode \code{what()} besitzen.
\begin{center}\small
\begin{tabular}{@{}ll@{}}
\toprule
\textbf{Kategorie} & \textbf{Beispiele} \\
\midrule
Logische Fehler / ung\"ultige Eingaben & \code{invalid\_argument}, \code{length\_error}, \code{logic\_error} \\
Laufzeitfehler & \code{runtime\_error}, \code{overflow\_error}, \code{out\_of\_range} \\
Speicher- \& Typfehler & \code{bad\_alloc}, \code{bad\_cast} \\
\bottomrule
\end{tabular}
\end{center}

\subsection{Eigene Exception-Klassen}
F\"ur fachlich spezifische Fehler -- erben von einer passenden Standard-Exception.
\begin{lstlisting}
class DivByZeroErr : public std::logic_error {
public:
    DivByZeroErr() : std::logic_error("Division durch 0") {}
};
// throw DivByZeroErr();
// catch (const DivByZeroErr& error) { std::cout << error.what(); }
\end{lstlisting}
\begin{tipp}[: Leitlinien]
\begin{itemize}
  \item Standard-Exceptions bevorzugen; eigene Klassen nur bei fachlich relevanten Fehlern.
  \item \textbf{Be specific} -- die spezifischste sinnvolle Exception werfen.
  \item \code{try}-Bl\"ocke klein halten; \code{catch} dort, wo man den Fehler versteht.
  \item Exceptions \textbf{als \code{const} Referenz} fangen (verhindert \emph{Object Slicing}).
  \item \code{catch (const std::exception\& e)} ist besser als \code{catch (...)}.
\end{itemize}
\end{tipp}

\subsection{RAII -- Resource Acquisition Is Initialization}
\begin{definition}[: RAII]
Eine Ressource (Speicher, Datei, DB-Verbindung, \dots) wird beim \textbf{Erzeugen} eines Objekts erworben und beim \textbf{Zerst\"oren} automatisch freigegeben -- Ressourcen sind an die Objektlebensdauer gebunden. Umgesetzt \"uber Konstruktor + Destruktor.
\end{definition}
\begin{lstlisting}
class File {
    FILE* f;
public:
    File(const std::string& name) { f = fopen(name.c_str(), "r"); }
    ~File() { if (f) fclose(f); }     // wird IMMER aufgerufen, auch bei Exception
};
void foo() {
    File f("test.txt");
    if (error) return;                // Datei wird trotzdem sauber geschlossen
}
\end{lstlisting}
Vorteile: keine Memory Leaks, Exception-Sicherheit, klarer Besitz, bessere Wartbarkeit. Voraussetzung: Objekt ist \textbf{lokal} (kein roher \code{new}-Pointer ohne Lebensdauerverwaltung).

\subsection{Smart Pointer}\label{sec:smartptr}
Rohe Pointer verursachen Memory Leaks, Dangling Pointer, Double-Free und unklare Besitzverh\"altnisse. \emph{Smart Pointer} (Wrapper-Klassen aus \code{<memory>}) verhalten sich wie Pointer (\"uberladen \code{*} und \code{->}), \"ubernehmen aber die Speicherverwaltung automatisch (RAII) und rufen selbst\"andig \code{delete} auf.
\begin{center}\small
\begin{tabular}{@{}lllll@{}}
\toprule
 & \textbf{Besitz} & \textbf{\#Besitzer} & \textbf{Z\"ahler} & \textbf{Zerst\"orung} \\
\midrule
\code{unique\_ptr} & exklusiv & genau 1 & nein & beim Verlassen des Besitzers \\
\code{shared\_ptr} & geteilt & beliebig & ja & wenn letzter \code{shared\_ptr} weg \\
\code{weak\_ptr} & kein Besitz & 0 & nutzt, erh\"oht nicht & nie \\
\bottomrule
\end{tabular}
\end{center}

\subsubsection*{\texttt{std::unique\_ptr} -- exklusiver Besitz}
\begin{lstlisting}
auto p1 = std::make_unique<int>(5);   // reserviert intern mit new
// auto p2 = p1;                      // FEHLER: Kopieren verboten
auto p2 = std::move(p1);              // Besitz uebertragen -> p1 == nullptr
\end{lstlisting}
Kein manuelles \code{delete} mehr (keine Double-Delete-Gefahr). Verwendung ohne Besitz\"ubergabe: \code{print(*p)} (Referenz) oder \code{foo(p.get())} (roher Pointer, Besitz bleibt bei \code{p}). Ideal f\"ur Kompositionen und \"uberall, wo fr\"uher \code{new} stand -- so effizient wie ein roher Pointer.

\subsubsection*{\texttt{std::shared\_ptr} -- geteilter Besitz}
\begin{lstlisting}
auto p1 = std::make_shared<int>(5);
auto p2 = p1;                  // beide besitzen; Zaehler = 2
std::cout << p1.use_count();   // Referenzzaehler abfragen
\end{lstlisting}
Das Objekt wird erst gel\"oscht, wenn der \textbf{letzte} Besitzer verschwindet (\emph{Referenzz\"ahler}). Langsamer als ein roher Pointer wegen des Z\"ahlers.

\subsubsection*{\texttt{std::weak\_ptr} -- l\"ost das Zyklusproblem}
Besitzen sich zwei \code{shared\_ptr} gegenseitig (A$\to$B, B$\to$A), bleibt der Z\"ahler stets $>0$ -- die Destruktoren laufen nie (Memory Leak). Ein \code{weak\_ptr} zeigt auf ein von \code{shared\_ptr} verwaltetes Objekt, \textbf{erh\"oht den Z\"ahler aber nicht}.
\begin{lstlisting}
std::shared_ptr<A> temp = objB->a.lock();  // weak_ptr nicht direkt dereferenzierbar
if (temp) { /* Objekt existiert noch */ }  // sonst temp == nullptr
\end{lstlisting}
\begin{merke}
Faustregel: \textbf{Wann immer eine Klasse/Funktion eine Ressource besitzt, Smart Pointer verwenden.} \code{unique\_ptr} als Standard (exklusiv), \code{shared\_ptr} nur bei echtem geteiltem Besitz, \code{weak\_ptr} zum Aufbrechen von Zyklen.
\end{merke}
\clearpage
